\documentclass{article}
\usepackage{graphicx} % Required for inserting images
\usepackage{ctex}
\usepackage{tikz}
\usetikzlibrary{calc,angles,quotes}
\usepackage{amsmath}

\newcommand{\p}{\partial}
\newcommand{\e}{\epsilon}
\newcommand{\sd}{\nabla\cdot}%散度
\newcommand{\xd}{\nabla \times}%旋度
\newcommand{\dv}[2]{\frac{\mathrm d #1}{\mathrm d #2}}
\newcommand{\dvv}[2]{\frac{\mathrm d^2 #1}{\mathrm d #2^2}}
\newcommand{\pdv}[2]{\frac{\partial #1}{\partial #2}}
\newcommand{\pdvv}[2]{\frac{\partial^2 #1}{\partial #2^2}}
\newcommand{\ti}{\times}
\newcommand{\cd}{\cdot}
\newcommand{\de}{\delta}
\title{电动力学第四章习题解答}
\author{wangjiahe060511}
\date{August 2026}

\begin{document}

\maketitle
\section{表面等离极化激元}
\subsection{(1)}
金属与介电质中的情况大致一样，我们仅验证一个
只需要验证
\begin{align}
    (\nabla ^2+k^2)E=0
\end{align}
即可。带入可得
\begin{align}
    k^2=k_1^2+k_3^{(1)2}
\end{align}
这是一个满足波动方程的电磁波。随后带入
\begin{align}
    \xd E=-\pdv{B}{t}
\end{align}
可得$B_x=B_z=0$,这代表这个波是TM的。
\subsection{}
\begin{align}
    S=\frac12 Re(E\times H^*)
\end{align}
带入可知只有x分量与z分量（因H只有y分量）
在介质侧，为
\begin{align}
    S_x^{(2)}&=\frac{\omega \e_0 \e_2 k_1}{2|k_3^{(2)}|^2}|E_0^2|e^{-2Im(k_3^{(2)z})}\\
    S_z^{(2)}
&=
\frac{\omega\epsilon_0\epsilon_2}{2}
|E_0|^2
e^{-2\mathrm{Im}(k_3^{(2)})z}
\frac{\mathrm{Re}(k_3^{(2)})}
{|k_3^{(2)}|^2}.
\end{align}
而在金属区，则有

\begin{align}
S_x^{(1)}
&=
\frac{\omega\epsilon_0\epsilon_1 k_1}
{2|k_3^{(1)}|^2}
|E_0|^2
e^{2\mathrm{Im}(k_3^{(1)})z},
\end{align}

以及

\begin{align}
S_z^{(1)}
&=
-\frac{\omega\epsilon_0\epsilon_1}{2}
|E_0|^2
e^{2\mathrm{Im}(k_3^{(1)})z}
\frac{\mathrm{Re}(k_3^{(1)})}
{|k_3^{(1)}|^2}.
\end{align}
无损耗，需要在两个区域能量均不沿z轴传播，也就是$$Re(k_3^{(1,2)})=0$$
\subsection{}
直接$$d=\frac{1}{Im(k_3^{(1)})}$$即可
\section{金属薄膜中的表面等离极化激元}

设金属薄膜的介电常数为$\epsilon_1$，厚度为$d$，
占据区域
\[
-\frac d2<z<\frac d2.
\]
薄膜下方介质的介电常数为$\epsilon_2$，
上方介质的介电常数为$\epsilon_3$。
不假设$\epsilon_2=\epsilon_3$。

由第一题已经知道，SPP为TM波，因此磁场只有$y$分量。
仍假设波沿$x$方向传播，并写成
\begin{align}
    H_y(x,z,t)=h(z)e^{ik_1x-i\omega t}.
\end{align}
在每一种均匀介质中，$H_y$均满足Helmholtz方程
\begin{align}
    \nabla^2 H_y+\epsilon_j k_0^2H_y=0,
    \qquad
    k_0=\frac{\omega}{c}.
\end{align}
代入上述形式可得
\begin{align}
    \dvv{h}{z}
    -
    \kappa_j^2 h=0,
\end{align}
其中定义
\begin{align}
    \kappa_j^2
    =
    k_1^2-\epsilon_jk_0^2.
\end{align}

为了得到束缚在金属薄膜附近的SPP，
上下两个无限大介质中的场必须在$|z|\rightarrow\infty$时衰减。
因此三区域中的解分别写为
\begin{align}
H_y^{(2)}
&=
A e^{\kappa_2(z+d/2)}
e^{ik_1x-i\omega t},
\qquad z<-\frac d2,
\\
H_y^{(1)}
&=
\left(
B e^{\kappa_1 z}
+
C e^{-\kappa_1 z}
\right)
e^{ik_1x-i\omega t},
\qquad
-\frac d2<z<\frac d2,
\\
H_y^{(3)}
&=
D e^{-\kappa_3(z-d/2)}
e^{ik_1x-i\omega t},
\qquad z>\frac d2.
\end{align}

注意金属薄膜的厚度有限，因此在金属内部不存在
$z\rightarrow\pm\infty$的边界条件，
故$e^{\kappa_1z}$和$e^{-\kappa_1z}$两项均必须保留。
这与第一题中半无限金属只有一个指数项不同。

由Maxwell方程
\begin{align}
    \xd H
    =
    \pdv{D}{t}
    =
    -i\omega\epsilon_0\epsilon_j E
\end{align}
可得
\begin{align}
    E_x^{(j)}
    =
    -\frac{i}{\omega\epsilon_0\epsilon_j}
    \pdv{H_y^{(j)}}{z}.
\end{align}
因此，在两个界面处需要满足
\begin{align}
    H_y\ \text{连续},
    \qquad
    \frac{1}{\epsilon}
    \pdv{H_y}{z}\ \text{连续}.
\end{align}

为了简化记号，定义
\begin{align}
    p_j=\frac{\kappa_j}{\epsilon_j}.
\end{align}

在下界面$z=-d/2$处，有
\begin{align}
    B e^{-\kappa_1d/2}
    +
    C e^{\kappa_1d/2}
    &=A,
\\
    p_1
    \left(
    B e^{-\kappa_1d/2}
    -
    C e^{\kappa_1d/2}
    \right)
    &=p_2A.
\end{align}

在上界面$z=d/2$处，有
\begin{align}
    B e^{\kappa_1d/2}
    +
    C e^{-\kappa_1d/2}
    &=D,
\\
    p_1
    \left(
    B e^{\kappa_1d/2}
    -
    C e^{-\kappa_1d/2}
    \right)
    &=-p_3D.
\end{align}

这四个方程共同决定$A,B,C,D$。
若希望存在非零的电磁场解，则方程组的系数行列式必须为零。
消去$A,B,C,D$后可得金属薄膜SPP的色散关系
\begin{align}
    \left(p_1^2+p_2p_3\right)
    \sinh(\kappa_1d)
    +
    p_1(p_2+p_3)
    \cosh(\kappa_1d)
    =0.
\end{align}
等价地，
\begin{align}
    \tanh(\kappa_1d)
    =
    -
    \frac{p_1(p_2+p_3)}
    {p_1^2+p_2p_3}.
\end{align}

这个方程体现了上下两个金属--介质界面之间的耦合。

当金属薄膜很厚，即
\begin{align}
    \kappa_1d\gg1,
\end{align}
有
\begin{align}
    \tanh(\kappa_1d)\rightarrow1.
\end{align}
于是色散关系变为
\begin{align}
    p_1^2+p_1p_2+p_1p_3+p_2p_3=0,
\end{align}
即
\begin{align}
    (p_1+p_2)(p_1+p_3)=0.
\end{align}
因此存在两种可能：
\begin{align}
    p_1+p_2=0,
\end{align}
或者
\begin{align}
    p_1+p_3=0.
\end{align}
它们分别正是金属与下方介质、金属与上方介质形成的
两个独立单界面SPP的色散关系。

因此，在$d$远大于SPP在金属中的穿透深度时，
上下两个界面上的SPP近似互不影响；
而当$d$与穿透深度可比时，
两个界面的倏逝场在金属内部发生重叠，
原来的两个独立SPP便发生耦合，形成整个三层结构的两个混合本征模。

特别地，如果上下介质恰好相同，即
\begin{align}
    \epsilon_2=\epsilon_3,
\end{align}
整个系统才具有关于$z=0$的镜面对称性。
此时两个混合模式可以进一步分成偶模和奇模，
金属内部的场分别可写为
\begin{align}
    H_y^{(+)}
    &\propto
    \cosh(\kappa_1z),
\\
    H_y^{(-)}
    &\propto
    \sinh(\kappa_1z).
\end{align}
\section{各向异性介质中电磁波的传播}

考虑均匀、非磁性的各向异性介质，
\begin{align}
    D_i=\epsilon_0\epsilon_{ij}E_j.
\end{align}
对于无自由电荷、无自由电流区域，Maxwell方程给出
\begin{align}
    \nabla\times\nabla\times E
    =
    k_0^2\boldsymbol{\epsilon}E,
    \qquad
    k_0=\frac{\omega}{c}.
\end{align}

对于平面波
\begin{align}
    E=E_0e^{i(\mathbf k\cdot\mathbf r-\omega t)},
\end{align}
有
\begin{align}
    \nabla\rightarrow i\mathbf k,
\end{align}
故
\begin{align}
    -\mathbf k\times(\mathbf k\times E)
    =
    k_0^2\boldsymbol{\epsilon}E.
\end{align}
利用
\begin{align}
    \mathbf k\times(\mathbf k\times E)
    =
    \mathbf k(\mathbf k\cdot E)-k^2E,
\end{align}
得到
\begin{align}
    \left(
    k^2I-\mathbf k\mathbf k^{T}
    -k_0^2\boldsymbol{\epsilon}
    \right)E=0.
\end{align}

即
\begin{align}
\begin{pmatrix}
k_y^2+k_z^2-k_0^2\epsilon_{11}
&
-k_xk_y-k_0^2\epsilon_{12}
&
-k_xk_z-k_0^2\epsilon_{13}
\\
-k_xk_y-k_0^2\epsilon_{21}
&
k_x^2+k_z^2-k_0^2\epsilon_{22}
&
-k_yk_z-k_0^2\epsilon_{23}
\\
-k_xk_z-k_0^2\epsilon_{31}
&
-k_yk_z-k_0^2\epsilon_{32}
&
k_x^2+k_y^2-k_0^2\epsilon_{33}
\end{pmatrix}
\begin{pmatrix}
E_x\\
E_y\\
E_z
\end{pmatrix}
=0.
\end{align}

为了存在非零电场，必须满足
\begin{align}
    \boxed{
    \det
    \left(
    k^2I-\mathbf k\mathbf k^{T}
    -k_0^2\boldsymbol{\epsilon}
    \right)=0
    }.
\end{align}

\subsection{介电张量的主轴}

对于无损耗的互易介质，
$\boldsymbol{\epsilon}$为实对称矩阵，因此可以取其三个本征矢
作为坐标轴，使其对角化：
\begin{align}
    \boldsymbol{\epsilon}
    =
    \begin{pmatrix}
    \epsilon_1&0&0\\
    0&\epsilon_2&0\\
    0&0&\epsilon_3
    \end{pmatrix}.
\end{align}

三个本征矢即介质的三个主轴，
三个本征值即三个主介电常数。

对于单轴晶体，取光轴为$z$轴，则
\begin{align}
    \boldsymbol{\epsilon}
    =
    \begin{pmatrix}
    \epsilon_\perp&0&0\\
    0&\epsilon_\perp&0\\
    0&0&\epsilon_\parallel
    \end{pmatrix}.
\end{align}

由于$x,y$方向完全等价，可以利用绕$z$轴的旋转对称性，
令波矢位于$xz$平面：
\begin{align}
    \mathbf k=(k_x,0,k_z).
\end{align}

此时本征方程化为
\begin{align}
\begin{pmatrix}
k_z^2-k_0^2\epsilon_\perp
&
0
&
-k_xk_z
\\
0
&
k_x^2+k_z^2-k_0^2\epsilon_\perp
&
0
\\
-k_xk_z
&
0
&
k_x^2-k_0^2\epsilon_\parallel
\end{pmatrix}
\begin{pmatrix}
E_x\\
E_y\\
E_z
\end{pmatrix}
=0.
\end{align}

行列式为零给出
\begin{align}
&
\left(
k_x^2+k_z^2-k_0^2\epsilon_\perp
\right)
\nonumber\\
&\qquad\times
\left[
\epsilon_\perp k_x^2
+\epsilon_\parallel k_z^2
-k_0^2\epsilon_\perp\epsilon_\parallel
\right]
=0.
\end{align}

因此存在两支本征模式。

\subsection{普通光与非常光}

第一支满足
\begin{align}
    \boxed{
    k_x^2+k_z^2
    =
    \epsilon_\perp k_0^2
    }.
\end{align}
这一支称为普通光（ordinary wave）。

由本征方程可知其电场为
\begin{align}
    \boxed{
    E_o=(0,E_y,0)
    }.
\end{align}

这里$y$方向并不是晶体本身特殊的方向，
而是由于我们已经把$\mathbf k$选在$xz$平面内。
一般而言，
\begin{align}
    \boxed{
    E_o\propto
    \hat{\mathbf a}\times\mathbf k
    },
\end{align}
其中$\hat{\mathbf a}$为光轴方向。

因此，对于给定的$\mathbf k$，
普通光只有一个独立的偏振本征方向。

第二支满足
\begin{align}
    \boxed{
    \frac{k_x^2}{\epsilon_\parallel}
    +
    \frac{k_z^2}{\epsilon_\perp}
    =
    k_0^2
    }.
\end{align}
这一支称为非常光（extraordinary wave）。

其电场位于$xz$平面内：
\begin{align}
    E_e=(E_x,0,E_z).
\end{align}

但$E_x,E_z$并不是两个独立的偏振自由度。
由
\begin{align}
    \nabla\cdot D=0
\end{align}
得到
\begin{align}
    \mathbf k\cdot D=0,
\end{align}
即
\begin{align}
    k_x\epsilon_\perp E_x
    +
    k_z\epsilon_\parallel E_z
    =0.
\end{align}
因此
\begin{align}
    \boxed{
    \frac{E_z}{E_x}
    =
    -\frac{\epsilon_\perp k_x}
    {\epsilon_\parallel k_z}
    }.
\end{align}

所以对于给定$\mathbf k$，非常光同样只有一个独立的偏振本征态。

普通光的等频面为球面，
而非常光的等频面为椭球面。
因此对于同一传播方向，一般存在两个不同的波数，
即产生双折射。

双折射的物理来源是介质沿光轴和平行于横向方向的
极化响应不同：
\begin{align}
    \epsilon_\parallel\neq\epsilon_\perp.
\end{align}
因此两种不同的本征偏振感受到不同的介质响应，
从而具有不同的色散关系。

\subsection{界面上的Fresnel问题}
\textbf{首先一个物理图像：oe光就是正常介质中的两种偏振！只不过在这种介质中，一个偏振有了偏移，不再垂直于波矢k}
下面考虑从各向同性介质入射到单轴晶体中的情况。

令界面为
\begin{align}
    z=0,
\end{align}
各向同性介质位于$z<0$，其相对介电常数为$\epsilon_i$；
单轴晶体位于$z>0$，光轴沿$z$方向。

令入射面为$xz$平面。

由于界面沿$x$方向具有平移对称性，
所有入射、反射以及透射波必须具有相同的切向波数：
\begin{align}
    \boxed{
    k_{ix}=k_{rx}=k_{ox}=k_{ex}\equiv k_x
    }.
\end{align}

入射介质中的法向波数为
\begin{align}
    k_{iz}
    =
    \sqrt{
    \epsilon_i k_0^2-k_x^2
    }.
\end{align}

需要特别注意的是，
在各向同性入射介质中不存在$o/e$光的区分，
而只有两种独立的偏振：
\begin{align}
    s:\quad E\perp\text{入射面},
\end{align}
以及
\begin{align}
    p:\quad E\parallel\text{入射面}.
\end{align}

在目前的几何构型中，恰好有
\begin{align}
    \boxed{
    s\longleftrightarrow o
    },
\end{align}
以及
\begin{align}
    \boxed{
    p\longleftrightarrow e
    }.
\end{align}

因此两种偏振完全解耦，
整个Fresnel问题可以分解为两个独立的$2\times2$边值问题：
\begin{align}
    \boxed{
    s_i\rightarrow s_r+o_t
    }
\end{align}
和
\begin{align}
    \boxed{
    p_i\rightarrow p_r+e_t
    }.
\end{align}

一般入射偏振
\begin{align}
    E_i=A_sE_s+A_pE_p
\end{align}
只需分别计算这两个问题，最后利用线性性叠加即可。

\subsection{$s$偏振入射：普通光}

对于$s$偏振，
\begin{align}
    E_i
    =
    (0,E_i,0),
\end{align}
反射波为
\begin{align}
    E_r
    =
    (0,E_r,0),
\end{align}
透射后的普通光为
\begin{align}
    E_o
    =
    (0,E_o,0).
\end{align}

普通光满足
\begin{align}
    k_x^2+k_{oz}^2
    =
    \epsilon_\perp k_0^2,
\end{align}
故
\begin{align}
    k_{oz}
    =
    \sqrt{
    \epsilon_\perp k_0^2-k_x^2
    }.
\end{align}

由
\begin{align}
    H
    =
    \frac{1}{\omega\mu_0}
    \mathbf k\times E
\end{align}
得到入射、反射及透射波的切向磁场：
\begin{align}
    H_{ix}
    &=
    -\frac{k_{iz}}{\omega\mu_0}E_i,
\\
    H_{rx}
    &=
    \frac{k_{iz}}{\omega\mu_0}E_r,
\\
    H_{ox}
    &=
    -\frac{k_{oz}}{\omega\mu_0}E_o.
\end{align}

界面上$E_y,H_x$连续，因此
\begin{align}
    E_i+E_r=E_o,
\end{align}
以及
\begin{align}
    k_{iz}(E_i-E_r)
    =
    k_{oz}E_o.
\end{align}

定义
\begin{align}
    r_o=\frac{E_r}{E_i},
    \qquad
    t_o=\frac{E_o}{E_i},
\end{align}
得到
\begin{align}
    \boxed{
    r_o
    =
    \frac{k_{iz}-k_{oz}}
    {k_{iz}+k_{oz}}
    },
\end{align}
以及
\begin{align}
    \boxed{
    t_o
    =
    \frac{2k_{iz}}
    {k_{iz}+k_{oz}}
    }.
\end{align}

因此普通光的Fresnel问题与普通各向同性介质中的
$s$偏振形式完全相同。

\subsection{$p$偏振入射：非常光}

对于$p$偏振，电场位于$xz$平面。

为了方便边界匹配，以切向电场$E_x$作为振幅。
在各向同性入射介质中有
\begin{align}
    \mathbf k_i\cdot E_i=0,
\end{align}
所以
\begin{align}
    E_{iz}
    =
    -\frac{k_x}{k_{iz}}E_{ix}.
\end{align}

由
\begin{align}
    H
    =
    \frac{1}{\omega\mu_0}
    \mathbf k\times E
\end{align}
得到
\begin{align}
    H_{iy}
    =
    \frac{\omega\epsilon_0\epsilon_i}
    {k_{iz}}E_{ix}.
\end{align}

对于反射$p$波，
\begin{align}
    H_{ry}
    =
    -\frac{\omega\epsilon_0\epsilon_i}
    {k_{iz}}E_{rx}.
\end{align}

非常光满足
\begin{align}
    \frac{k_x^2}{\epsilon_\parallel}
    +
    \frac{k_{ez}^2}{\epsilon_\perp}
    =
    k_0^2,
\end{align}
故
\begin{align}
    \boxed{
    k_{ez}
    =
    \sqrt{
    \epsilon_\perp
    \left(
    k_0^2
    -
    \frac{k_x^2}{\epsilon_\parallel}
    \right)
    }
    }.
\end{align}

非常光具有
\begin{align}
    E_e=(E_{ex},0,E_{ez}),
    \qquad
    H_e=(0,H_{ey},0).
\end{align}

由
\begin{align}
    \mathbf k\times H
    =
    -\omega D
\end{align}
的$x$分量可得
\begin{align}
    k_{ez}H_{ey}
    =
    \omega\epsilon_0\epsilon_\perp E_{ex}.
\end{align}
因此
\begin{align}
    H_{ey}
    =
    \frac{\omega\epsilon_0\epsilon_\perp}
    {k_{ez}}
    E_{ex}.
\end{align}

界面上$E_x,H_y$连续，故
\begin{align}
    E_{ix}+E_{rx}=E_{ex},
\end{align}
以及
\begin{align}
    \frac{\epsilon_i}{k_{iz}}
    (E_{ix}-E_{rx})
    =
    \frac{\epsilon_\perp}{k_{ez}}
    E_{ex}.
\end{align}

定义
\begin{align}
    r_e=\frac{E_{rx}}{E_{ix}},
    \qquad
    t_e=\frac{E_{ex}}{E_{ix}},
\end{align}
得到
\begin{align}
    \boxed{
    r_e
    =
    \frac{
    \epsilon_i k_{ez}
    -
    \epsilon_\perp k_{iz}
    }{
    \epsilon_i k_{ez}
    +
    \epsilon_\perp k_{iz}
    }
    },
\end{align}
以及
\begin{align}
    \boxed{
    t_e
    =
    \frac{
    2\epsilon_i k_{ez}
    }{
    \epsilon_i k_{ez}
    +
    \epsilon_\perp k_{iz}
    }
    }.
\end{align}

因此当前几何下的Fresnel问题最终可总结为
\begin{align}
    \boxed{
    s\rightarrow o,
    \qquad
    p\rightarrow e
    }.
\end{align}

任意入射偏振只需分解为$s,p$两部分，
分别应用上述两套Fresnel系数后再叠加即可。

\subsection{两种模式的能流}

时间平均Poynting矢量为
\begin{align}
    S
    =
    \frac12
    \operatorname{Re}
    \left(
    E\times H^*
    \right).
\end{align}

对于普通光，
\begin{align}
    E_o=(0,E_y,0),
\end{align}
而
\begin{align}
    H_o
    =
    \frac{E_y}{\omega\mu_0}
    (-k_z,0,k_x).
\end{align}

因此
\begin{align}
    \boxed{
    S_o
    =
    \frac{|E_y|^2}
    {2\omega\mu_0}
    (k_x,0,k_z)
    }.
\end{align}
即
\begin{align}
    \boxed{
    S_o\parallel k_o
    }.
\end{align}

普通光的相位传播方向与能量传播方向一致。

对于非常光，
\begin{align}
    E_e=(E_x,0,E_z),
    \qquad
    H_e=(0,H_y,0).
\end{align}

由
\begin{align}
    \mathbf k\times H=-\omega D
\end{align}
得到
\begin{align}
    D_x
    =
    \frac{k_z}{\omega}H_y,
    \qquad
    D_z
    =
    -\frac{k_x}{\omega}H_y.
\end{align}

又因为
\begin{align}
    D_x
    =
    \epsilon_0\epsilon_\perp E_x,
    \qquad
    D_z
    =
    \epsilon_0\epsilon_\parallel E_z,
\end{align}
所以
\begin{align}
    E_x
    =
    \frac{k_z}
    {\omega\epsilon_0\epsilon_\perp}H_y,
\end{align}
以及
\begin{align}
    E_z
    =
    -\frac{k_x}
    {\omega\epsilon_0\epsilon_\parallel}H_y.
\end{align}

于是
\begin{align}
    \boxed{
    S_e
    =
    \frac{|H_y|^2}
    {2\omega\epsilon_0}
    \left(
    \frac{k_x}{\epsilon_\parallel},
    0,
    \frac{k_z}{\epsilon_\perp}
    \right)
    }.
\end{align}

因此一般有
\begin{align}
    \boxed{
    S_e\not\parallel k_e
    }.
\end{align}

非常光的相位传播方向与实际能量传播方向不同。

事实上，非常光的等频面为
\begin{align}
    F(\mathbf k)
    =
    \frac{k_x^2}{\epsilon_\parallel}
    +
    \frac{k_z^2}{\epsilon_\perp}
    -
    k_0^2
    =0.
\end{align}
有
\begin{align}
    \nabla_{\mathbf k}F
    =
    2
    \left(
    \frac{k_x}{\epsilon_\parallel},
    0,
    \frac{k_z}{\epsilon_\perp}
    \right),
\end{align}
故
\begin{align}
    \boxed{
    S_e\parallel\nabla_{\mathbf k}F
    }.
\end{align}

即非常光的能流沿等频椭圆的法线方向，
而不沿波矢本身。

\subsection{反射率与透射率}

Fresnel系数$r,t$是场振幅系数，
真正的能量反射率和透射率应由Poynting矢量计算。

对于普通光，
\begin{align}
    R_o=|r_o|^2,
\end{align}
并有
\begin{align}
    \boxed{
    T_o
    =
    \frac{k_{oz}}{k_{iz}}
    |t_o|^2
    }.
\end{align}

对于非常光，定义
\begin{align}
    Y_i^{(p)}
    =
    \frac{\omega\epsilon_0\epsilon_i}{k_{iz}},
\end{align}
以及
\begin{align}
    Y_e
    =
    \frac{\omega\epsilon_0\epsilon_\perp}{k_{ez}},
\end{align}
则
\begin{align}
    \boxed{
    T_e
    =
    \frac{Y_e}{Y_i^{(p)}}
    |t_e|^2
    }.
\end{align}

对于无损介质，应满足能量守恒
\begin{align}
    \boxed{
    R_o+T_o=1,
    \qquad
    R_e+T_e=1
    }.
\end{align}

\subsection{如果折射率变成复数}

若介质具有损耗，
介电常数必须写成复数。
对于单轴晶体，
\begin{align}
    \epsilon_\perp
    =
    \epsilon_\perp'
    +
    i\epsilon_\perp'',
\end{align}
以及
\begin{align}
    \epsilon_\parallel
    =
    \epsilon_\parallel'
    +
    i\epsilon_\parallel''.
\end{align}

此时色散方程形式不变，
但是波矢的解一般成为复数。

对于实际的界面问题，
切向波数$k_x$由入射波确定，
因此主要是两个法向波数$k_z$变为复数。

普通光有
\begin{align}
    \boxed{
    k_{oz}
    =
    \sqrt{
    \epsilon_\perp k_0^2-k_x^2
    }
    }.
\end{align}

非常光有
\begin{align}
    \boxed{
    k_{ez}
    =
    \sqrt{
    \epsilon_\perp
    \left(
    k_0^2
    -
    \frac{k_x^2}{\epsilon_\parallel}
    \right)
    }
    }.
\end{align}

由于$\epsilon_\perp,\epsilon_\parallel$为复数，
一般有
\begin{align}
    k_{oz}
    =
    k_{oz}'
    +
    ik_{oz}'',
\end{align}
以及
\begin{align}
    k_{ez}
    =
    k_{ez}'
    +
    ik_{ez}''.
\end{align}

对应的空间因子分别为
\begin{align}
    e^{ik_{oz}z}
    =
    e^{ik_{oz}'z}
    e^{-k_{oz}''z},
\end{align}
以及
\begin{align}
    e^{ik_{ez}z}
    =
    e^{ik_{ez}'z}
    e^{-k_{ez}''z}.
\end{align}

因此两个本征模式的场振幅衰减率分别为
\begin{align}
    \boxed{
    \alpha_o
    =
    \operatorname{Im}(k_{oz})
    },
\end{align}
以及
\begin{align}
    \boxed{
    \alpha_e
    =
    \operatorname{Im}(k_{ez})
    }.
\end{align}

一般情况下
\begin{align}
    \boxed{
    \alpha_o\neq\alpha_e
    }.
\end{align}

普通光只感受到$\epsilon_\perp$方向的介质响应，
而非常光同时包含垂直与平行于光轴的电场分量，
因此其衰减同时受到
$\epsilon_\perp''$和$\epsilon_\parallel''$的影响。

更一般地，若介电常数为复矩阵
\begin{align}
    \boldsymbol{\epsilon}
    =
    \boldsymbol{\epsilon}'
    +
    i\boldsymbol{\epsilon}'',
\end{align}
则仍然通过
\begin{align}
    \boxed{
    \det
    \left[
    k^2I-\mathbf k\mathbf k^T
    -k_0^2\boldsymbol{\epsilon}
    \right]
    =0
    }
\end{align}
求取允许的复波矢。

介电常数矩阵的虚部最终通过这一复本征方程
转化为两支模式波矢的虚部，从而决定各自的衰减率。
\section{电磁场的角动量和偏振的关系}
\subsection{}
从$$L=\frac{1}{\mu_0 c^2}\int r\ti (E\ti B)dV$$开始推导
$$L=\frac{1}{\mu_0 c^2}\int E\ti A+E_j (r\ti \nabla)A_jdV$$ \\
首先带入$B=\xd A$,然后逐层的把矢量运算转化成指标运算，最后积分内部可变为
\begin{align}
    \e_{ijk}r_j\e_{klm}E_l\e_{mno}\p_nA_o
\end{align}
然后利用$$\e_{ijk}\e_{ilm}=\de_{jl}\de_{km}-\de_{jm}\de_{kl}$$一通化简之后得到
\begin{align}
\e_{ijk}r_j(E_l\p_k A_l-E_l\p_l A_k)
\end{align}
对减号前后分别分部积分，注意到所有的全微分($\p \text{然后一套}$)在无穷远边界上足够快的衰减为零（此处用了高斯定理）
就把第一部分化成\begin{align}
    A_l\e_{ijk}\p_k(r_jE_l)
\end{align}
注意所有的Levi-civita符号都可以任意提出。里面乘积求导写开，结合$\p_k r_j=\de_{kj}$
就是\begin{align}
    E_l \e_{ijk}r_j\p_k A_l
\end{align}
写回向量运算模式也就是
\begin{align}
    E_l(r\ti \nabla)A_l
\end{align}
完全一样的，减号右侧就是$-E\ti A$,于是就得到了整体的形式。
\subsection{}
选取规范令$E=-\pdv{A}{t}$,有
\begin{align}
A=\frac{1}{8\pi^3}\int \e_\lambda a_\lambda e^{ikx-iwt}+\e^*_\lambda a^*_\lambda e^{-ikx+iwt}
\end{align}
叉乘起来，注意两者不用同一套k积分，最终是六重，然后忽略所有的时间震荡项，有
\[
\begin{aligned}
\mathbf E\times\mathbf A
=i\int\frac{d^3k\,d^3k'}{(2\pi)^6}\omega\Big[
+&(\mathbf F\times\mathbf F'^*)e^{i(\theta-\theta')}\\
-&(\mathbf F^*\times\mathbf F')e^{-i(\theta-\theta')}\\
\Big]
\end{aligned}
\]
其中\[
\mathbf F(\mathbf k)=\sum_{\lambda}\boldsymbol\epsilon_\lambda(\mathbf k)a_\lambda(\mathbf k),
\qquad
\theta=\mathbf k\cdot\mathbf x-\omega t
\]
整体的计算4.133中的对空间积分，所有的指数变成Dirac函数。
最终变成\[
\boxed{
\mathbf L_{\rm spin}
=\frac{2i}{\mu_0c^2}
\int\frac{d^3k}{(2\pi)^3}\omega
(\mathbf F\times\mathbf F^*)
}
\]分别带入\[
\mathbf e_\pm=\frac{\mathbf e_1\pm i\mathbf e_2}{\sqrt2},
\qquad
\mathbf e_1\times\mathbf e_2=\hat{\mathbf k}
\]\[
\mathbf F=a_+\mathbf e_++a_-\mathbf e_-
\]即可证出。
\subsubsection{电磁场的总能量}

真空中电磁场的总能量为
\begin{align}
    W=\frac12\int d^3x
    \left(
    \epsilon_0 E^2+\frac{1}{\mu_0}B^2
    \right).
\end{align}

沿用前面的记号
\begin{align}
    \mathbf F(\mathbf k)
    &=\sum_\lambda \boldsymbol\epsilon_\lambda a_\lambda,\\
    \mathbf G(\mathbf k)
    &=\mathbf k\times\mathbf F(\mathbf k),\\
    \theta&=\mathbf k\cdot\mathbf x-\omega t.
\end{align}

电场与磁场分别为
\begin{align}
    \mathbf E
    &=i\int\frac{d^3k}{(2\pi)^3}\omega
    \left(
    \mathbf F e^{i\theta}
    -\mathbf F^*e^{-i\theta}
    \right),\\
    \mathbf B
    &=i\int\frac{d^3k}{(2\pi)^3}
    \left(
    \mathbf G e^{i\theta}
    -\mathbf G^*e^{-i\theta}
    \right).
\end{align}

首先计算电场能量。将$\mathbf E$点乘自身，
并舍去时间平均后为零的振荡项，得到
\begin{align}
    \overline{E^2}
    =\int\frac{d^3k\,d^3k'}{(2\pi)^6}
    \omega\omega'
    \Big[
    &\mathbf F\cdot\mathbf F'^*
    e^{i(\theta-\theta')}\\
    +&\mathbf F^*\cdot\mathbf F'
    e^{-i(\theta-\theta')}
    \Big].
\end{align}

随后利用空间积分
\begin{align}
    \int d^3x\,
    e^{i(\mathbf k-\mathbf k')\cdot\mathbf x}
    =(2\pi)^3\delta^{(3)}(\mathbf k-\mathbf k'),
\end{align}
消去$\mathbf k'$积分，两项相同，因此
\begin{align}
    \boxed{
    \int d^3x\,\overline{E^2}
    =2\int\frac{d^3k}{(2\pi)^3}
    \omega^2|\mathbf F|^2
    }.
\end{align}

磁场部分完全类似，只需将
$\omega\mathbf F$替换为$\mathbf G$，得到
\begin{align}
    \boxed{
    \int d^3x\,\overline{B^2}
    =2\int\frac{d^3k}{(2\pi)^3}
    |\mathbf k\times\mathbf F|^2
    }.
\end{align}

由于电磁波满足横波条件
\begin{align}
    \mathbf k\cdot\mathbf F=0,
\end{align}
因此
\begin{align}
    |\mathbf k\times\mathbf F|^2
    &=k^2|\mathbf F|^2
    -|\mathbf k\cdot\mathbf F|^2\\
    &=k^2|\mathbf F|^2.
\end{align}

再利用真空色散关系
\begin{align}
    \omega=ck,\qquad
    \frac{1}{\mu_0}=\epsilon_0c^2,
\end{align}
可知电场与磁场的能量贡献相等，即
\begin{align}
    \epsilon_0\int d^3x\,\overline{E^2}
    =\frac{1}{\mu_0}
    \int d^3x\,\overline{B^2}.
\end{align}

所以总能量为
\begin{align}
    \overline W
    =\frac{2}{\mu_0c^2}
    \int\frac{d^3k}{(2\pi)^3}
    \omega^2|\mathbf F(\mathbf k)|^2.
\end{align}

最后代入圆偏振基
\begin{align}
    \mathbf F=a_+\mathbf e_++a_-\mathbf e_-.
\end{align}

利用正交归一关系
\begin{align}
    \mathbf e_\lambda^*
    \cdot\mathbf e_{\lambda'}
    =\delta_{\lambda\lambda'},
\end{align}
得到
\begin{align}
    |\mathbf F|^2
    &=\mathbf F^*\cdot\mathbf F\\
    &=|a_+|^2+|a_-|^2.
\end{align}

最终
\begin{align}
    \boxed{
    \overline W
    =\frac{2}{\mu_0c^2}
    \int\frac{d^3k}{(2\pi)^3}
    \omega^2
    \left(
    |a_+(\mathbf k)|^2
    +|a_-(\mathbf k)|^2
    \right)
    }.
\end{align}

\section{菲涅尔公式的另外形式}
题意不很明确俺觉得，带入折射定律消去俩折射率就行（由ChatGPT 6Sol建议）
垂直偏振（正弦公式）
\[
\boxed{
\begin{aligned}
r_\perp&=-\frac{\sin(i-r)}{\sin(i+r)}\\[6pt]
t_\perp&=\frac{2\sin r\cos i}{\sin(i+r)}
\end{aligned}
}
\]


平行偏振（正切公式）
\[
\boxed{
\begin{aligned}
r_\parallel&=\frac{\tan(i-r)}{\tan(i+r)}\\[6pt]
t_\parallel&=\frac{2\sin r\cos i}
{\sin(i+r)\cos(i-r)}
\end{aligned}
}
\]
\section{均匀磁场中电磁波的传播}

利用
\begin{align}
    \epsilon_{ij}=\epsilon_0(\delta_{ij}+\chi_{ij})
\end{align}
求出介电常数张量，再分析其本征值、本征偏振以及不同传播方向下的色散关系。

\subsection{介电常数张量及其本征值}

由于静磁场均匀，可以选择坐标系使其沿 $z$ 轴：
\begin{align}
    \mathbf B_0=B_0\hat{\mathbf z},\qquad
    \mathbf b=(0,0,1).
\end{align}
将其代入题目给出的 $\chi_{ij}$，得到
\begin{align}
    \boldsymbol\chi
    =\begin{pmatrix}
    -\dfrac{\omega_p^2}{\omega^2-\omega_B^2}
    &i\dfrac{\omega_p^2\omega_B}{\omega(\omega^2-\omega_B^2)}&0\\[8pt]
    -i\dfrac{\omega_p^2\omega_B}{\omega(\omega^2-\omega_B^2)}
    &-\dfrac{\omega_p^2}{\omega^2-\omega_B^2}&0\\[8pt]
    0&0&-\dfrac{\omega_p^2}{\omega^2}
    \end{pmatrix}.
\end{align}
因此相对介电常数张量为
\begin{align}
    \boldsymbol\epsilon_r
    \equiv\frac{\boldsymbol\epsilon}{\epsilon_0}
    =\mathbf I+\boldsymbol\chi
    =\begin{pmatrix}
    S&iD&0\\
    -iD&S&0\\
    0&0&P
    \end{pmatrix},
\end{align}
其中
\begin{align}
    S&=1-\frac{\omega_p^2}{\omega^2-\omega_B^2},\\
    D&=\frac{\omega_p^2\omega_B}{\omega(\omega^2-\omega_B^2)},\\
    P&=1-\frac{\omega_p^2}{\omega^2}.
\end{align}
三个本征值为
\begin{align}
    \boxed{\epsilon_+=\epsilon_0(S+D),\qquad
    \epsilon_-=\epsilon_0(S-D),\qquad
    \epsilon_z=\epsilon_0P.}
\end{align}

对 $\lambda_+=S+D$，本征方程给出 $E_y=-iE_x$；对 $\lambda_-=S-D$，则有 $E_y=iE_x$。归一化后可取三个本征矢为
\begin{align}
    \boxed{
    \mathbf e_+=\frac{1}{\sqrt2}\begin{pmatrix}1\\-i\\0\end{pmatrix},\qquad
    \mathbf e_- =\frac{1}{\sqrt2}\begin{pmatrix}1\\i\\0\end{pmatrix},\qquad
    \mathbf e_z=\begin{pmatrix}0\\0\\1\end{pmatrix}.}
\end{align}
前两个本征矢是复矢量，这意味着不再是之前处理的线偏振，而必然是圆偏振，因为有些落后了i的相位。
\\
其他的讨论，请见双折射一节。
\\

\section{波包传播与不确定关系}

由式（4.56）和（4.61），

\begin{align}
u(x,t)
&=\frac{1}{\sqrt{2\pi}}
\int_{-\infty}^{\infty}
A(k)e^{i[kx-\omega(k)t]}\,dk\\
&=e^{i(k_0v_g-\omega_0)t}
\frac{1}{\sqrt{2\pi}}
\int_{-\infty}^{\infty}
A(k)e^{ik(x-v_gt)}\,dk\\
&=e^{i(k_0v_g-\omega_0)t}u(x-v_gt,0).
\end{align}

即验证式（4.62）。

\section{波导中的横向和纵向电磁场的关系}

令电磁场具有因子 $e^{ik_zz-i\omega t}$，由式（4.83）和（4.84），

\begin{align}
k_z\mathbf E_\perp
+\omega\mathbf e_3\times\mathbf B_\perp
&=-i\nabla_\perp E_z,\\
k_z\mathbf B_\perp
-\mu\epsilon\omega\mathbf e_3\times\mathbf E_\perp
&=-i\nabla_\perp B_z.
\end{align}

利用
\[
\mathbf e_3\times(\mathbf e_3\times\mathbf A)
=-\mathbf A,
\qquad
k_\perp^2=\mu\epsilon\omega^2-k_z^2,
\]
消去另一横向场分量，得到

\begin{align}
\boxed{
\mathbf E_\perp
=\frac{i}{k_\perp^2}
\left[
k_z\nabla_\perp E_z
-\omega\mathbf e_3\times\nabla_\perp B_z
\right]
},\\[6pt]
\boxed{
\mathbf B_\perp
=\frac{i}{k_\perp^2}
\left[
k_z\nabla_\perp B_z
+\mu\epsilon\omega\mathbf e_3\times\nabla_\perp E_z
\right]
}.
\end{align}

即验证式（4.87）。

\section{金属波导单位长度的损耗}

由金属中的电磁场方程（4.39），良导体的趋肤深度为

\[
\delta=\sqrt{\frac{2}{\mu\sigma\omega}},
\]

电场沿金属内部法向深度 $s$ 衰减：

\[
\mathbf E(s)=\mathbf E(0)e^{-(1-i)s/\delta},
\qquad
|\mathbf E(0)|^2
=\frac{2}{\sigma^2\delta^2}
|\mathbf n\times\mathbf H|^2.
\]

单位面积上的平均焦耳损耗为

\begin{align}
p_{\rm loss}
&=\frac{\sigma}{2}\int_0^\infty
|\mathbf E(s)|^2\,ds\\
&=\frac{\sigma\delta}{4}|\mathbf E(0)|^2\\
&=\frac{1}{2\sigma\delta}
|\mathbf n\times\mathbf H|^2.
\end{align}

沿波导横截面边界 $C$ 积分，得到

\[
\boxed{
-\frac{dP}{dz}
=\frac{1}{2\sigma\delta}
\oint_C|\mathbf n\times\mathbf H|^2\,dl
}.
\]

即验证式（4.110）。

\section{声明}
本人推导学习时，忘记及时整理7,8,9题，故该三题答案由GPT生成，其余题目由本人撰写，难免有错误，欢迎指出！
\end{document}